\chapter{Distributions, Expectation, and Variance}
\label{chap:distributions}

\section{Exponential waiting times}
\label{sec:exponential}

Imagine watching a help desk's inbox. Start a stopwatch now and stop it
when the next message arrives. The number on the stopwatch is \(T\), the
waiting time in minutes. On one occasion we might wait \(0.3\) minutes;
on another, \(7\) minutes. We want to answer questions such as: what is
the probability that we are still waiting after five minutes?

\subsection{A model of arrivals without a schedule}

First build a simple model using one-second intervals. Suppose that in
each interval a message arrives with probability \(1/300\), independently
of what happened in earlier intervals. Most seconds contain no message.
Occasionally two successful seconds occur close together; occasionally
there is a long gap.

In this model, waiting more than five minutes means getting no message in
any of \(300\) consecutive seconds. By independence,
\[
  \Prob(T>5)=\left(1-\frac1{300}\right)^{300}\approx0.3673.
\]
A long wait is therefore quite plausible. The model does not place an
arrival at regular five-minute intervals.

Now shorten the time intervals. Write \(\lambda\) for the arrival rate,
measured in messages per minute. A small interval of length \(\Delta\)
minutes has arrival probability \(\lambda\Delta\).
In our example, \(\lambda=0.2\), because
\(0.2\times(1/60)=1/300\).
The rate describes how frequently messages arrive on average; it is not
itself a probability.

Divide a time span of \(t\) minutes into \(n\) equal intervals. In the
discrete approximation, each has arrival probability \(\lambda t/n\).
The probability that all of them are empty is
\[
  \left(1-\frac{\lambda t}{n}\right)^n.
\]
Letting the intervals become arbitrarily short gives
\begin{equation}
  \Prob(T>t)
  =\lim_{n\to\infty}\left(1-\frac{\lambda t}{n}\right)^n
  =e^{-\lambda t},\qquad t\ge0.
  \label{eq:exponential-tail}
\end{equation}
This is the probability that \emph{no message has arrived by time \(t\)}.
The exponential function appears because we multiply the no-arrival
probabilities of many short intervals. The resulting continuous waiting
time has an \emph{exponential distribution}.

The assumptions describe an inbox with no timetable and no change in
activity during the period being studied. They would be unsuitable for
notifications sent every hour or an inbox that becomes much busier at
lunchtime.

\subsection{From the waiting probability to the density}

A message has either arrived by time \(t\), or it has not. Thus the CDF is
\[
  F_T(t)=\Prob(T\le t)=1-e^{-\lambda t},\qquad t\ge0.
\]
Differentiating this CDF, using the relation from Chapter~1, gives the
density of the waiting time.

\begin{definition}[Exponential distribution]
For \(\lambda>0\), the notation \(T\sim\operatorname{Exp}(\lambda)\) means
\[
  f_T(t)=
  \begin{cases}
    \lambda e^{-\lambda t},&t\ge0,\\
    0,&t<0.
  \end{cases}
\]
The parameter \(\lambda\) is the arrival rate.
\end{definition}

Increasing \(\lambda\) makes short waits more likely. For example, at rate
\(0.2\) per minute the probability of waiting more than five minutes is
\(e^{-1}\approx0.3679\). At rate \(0.5\), that probability drops to
\(e^{-2.5}\approx0.0821\). Figure~\ref{fig:exponential-wait} shows this
comparison directly.

\begin{figure}[htbp]
  \centering
  \begin{tikzpicture}
\begin{axis}[
  width=0.94\textwidth,height=6.3cm,
  axis lines=left,xmin=0,xmax=20,ymin=0,ymax=1.05,
  xlabel={Time waited \(t\) (minutes)},
  ylabel={Probability of still waiting},
  tick label style={font=\small},label style={font=\small},
  legend style={draw=none,font=\small,at={(0.98,0.98)},anchor=north east},
  grid=major,grid style={black!8},domain=0:20,samples=150,
]
  \addplot[CourseBlue,very thick] {exp(-0.2*x)};
  \addlegendentry{\(\lambda=0.2\) messages/minute}
  \addplot[orange!80!black,thick,dashed] {exp(-0.5*x)};
  \addlegendentry{\(\lambda=0.5\) messages/minute}
  \addplot[teal!80!black,thick,dashdotted] {exp(-0.1*x)};
  \addlegendentry{\(\lambda=0.1\) messages/minute}
  \addplot[black!40,densely dashed,forget plot] coordinates {(5,0) (5,1)};
  \addplot[CourseBlue,only marks,mark=*,mark size=2pt,forget plot]
    coordinates {(5,0.367879)};
\end{axis}
\end{tikzpicture}

  \caption{The probability that the inbox is still empty after \(t\) minutes.
  Every curve starts at one: at time zero we have not yet waited.
  A busier inbox has a larger rate and a faster decrease in the probability
  of still waiting. The dashed vertical line marks five minutes.}
  \label{fig:exponential-wait}
\end{figure}

For an interval question, subtract two waiting probabilities. With
\(\lambda=0.2\),
\[
  \Prob(2<T\le5)
  =\Prob(T>2)-\Prob(T>5)
  =e^{-0.4}-e^{-1}\approx0.3024.
\]
We keep the cases that are still waiting at minute two and remove those
that are still waiting at minute five.

A different practical question is how long we must wait to have seen a
message in \(95\%\) of repetitions. Solve
\[
  1-e^{-0.2t}=0.95
  \quad\Longrightarrow\quad
  t=-\frac{\log(0.05)}{0.2}\approx14.98\text{ minutes}.
\]
This cutoff is the \(95\)-th percentile of the waiting time. More
generally, the \(q\)-quantile is
\(t_q=-\log(1-q)/\lambda\), for \(0<q<1\).

\subsection{Does a long wait make an arrival overdue?}

Suppose we have already waited ten minutes with no message. In the model,
the next short interval has the same arrival chance as it did at the
start. Earlier empty intervals do not change that chance. Therefore the
probability of waiting another five minutes is still
\(e^{-0.2\cdot5}\approx0.3679\).

We can verify this using conditional probability. For \(s,t\ge0\),
\[
  \Prob(T>s+t\mid T>s)
  =\frac{\Prob(T>s+t)}{\Prob(T>s)}
  =\frac{e^{-\lambda(s+t)}}{e^{-\lambda s}}
  =e^{-\lambda t}.
\]
The right-hand side is \(\Prob(T>t)\). The distribution of the remaining
wait is the same as the distribution of the original wait. This property
is called \emph{memorylessness}.

A scheduled bus behaves differently: if it comes every ten minutes,
waiting nine minutes tells us something about the remaining wait.
The exponential model is appropriate only when this kind of schedule,
ageing, or change in activity is absent.

\section{Expectation and functions of a random variable}
\label{sec:expectation}

The probability of a long wait answers one question. The average amount
of waiting answers another. If we repeat the inbox experiment independently
many times and average the recorded times, what value should that average
be close to?

Start with a simpler example. Suppose a task takes one minute with
probability \(3/4\) and nine minutes with probability \(1/4\).
In many repetitions, about three quarters of the tasks take one minute
and one quarter take nine minutes. Their average is close to
\[
  \frac34\cdot1+\frac14\cdot9=3\text{ minutes}.
\]
This probability-weighted average is the \emph{expectation}. It can be
three minutes even though no individual task takes three minutes.

\subsection{Averaging with probabilities}

For a discrete random variable, weight each possible value by its
probability and add. For a continuous variable, a small interval near
\(x\) has probability approximately \(f_X(x)\Delta x\), so its contribution
to the average is approximately \(x f_X(x)\Delta x\).
Summing these contributions over increasingly small intervals gives an
integral.

\begin{definition}[Expectation]
The expectation of \(X\) is
\[
  \Expect[X]=
  \begin{cases}
    \displaystyle\sum_x x\,p_X(x),&X\text{ discrete},\\[6pt]
    \displaystyle\int_{-\infty}^{\infty}x f_X(x)\,dx,
      &X\text{ has a density}.
  \end{cases}
\]
For a finite expectation we require \(\Expect[|X|]<\infty\): the sum or
integral of absolute contributions must converge.
\end{definition}

This condition matters when very large values, although rare, contribute
so much that no finite average exists. All the distributions in the main
part of this chapter have finite means. The Student extension provides
an example where the condition can fail.

For exponential waiting times, the definition gives
\[
  \Expect[T]=\int_0^\infty t\lambda e^{-\lambda t}\,dt.
\]
To evaluate it, integrate by parts with \(u=t\) and
\(dv=\lambda e^{-\lambda t}\,dt\):
\[
  \Expect[T]
  =[-t e^{-\lambda t}]_0^\infty
    +\int_0^\infty e^{-\lambda t}\,dt
  =\frac1\lambda.
\]
At rate \(0.2\) messages per minute, the mean wait is five minutes.
This is an average over repeated waits, not an appointment for the next
arrival. We already found that more than a third of waits exceed five
minutes.

\subsection{First transform the observation, then average}

Sometimes we want to average something other than the waiting time
itself. For example, assign a response score
\[
  g(t)=e^{-at},\qquad a>0.
\]
An immediate response scores one; longer waits receive smaller scores.
The parameter \(a\) controls how quickly the score decreases with time.

To find the average score, calculate the score for each possible wait and
then weight it by the probability of that wait:
\[
  \Expect[g(T)]=\int_0^\infty g(t)f_T(t)\,dt.
\]
For this particular score,
\[
  \Expect[e^{-aT}]
  =\int_0^\infty e^{-at}\lambda e^{-\lambda t}\,dt
  =\frac{\lambda}{\lambda+a}.
\]
This calculation illustrates the general rule
\begin{equation}
  \Expect[g(X)]=
  \begin{cases}
    \displaystyle\sum_x g(x)p_X(x),&X\text{ discrete},\\[6pt]
    \displaystyle\int_{-\infty}^{\infty}g(x)f_X(x)\,dx,
      &X\text{ has a density}.
  \end{cases}
  \label{eq:expectation-function}
\end{equation}
We assume \(\Expect[|g(X)|]<\infty\). The rule lets us use the distribution
of \(X\) directly, without first deriving the distribution of \(g(X)\).

Scoring the average wait is a different operation. For \(a=\lambda=1\),
\[
  \Expect[e^{-T}]=\frac12,
  \qquad
  e^{-\Expect[T]}=e^{-1}\approx0.3679.
\]
The nonlinear transformation makes the order of the two operations matter.

Another useful choice is \(g(t)=t^2\). Large waits receive much more weight
after squaring. Again integrating by parts,
\[
  \Expect[T^2]
  =\int_0^\infty t^2\lambda e^{-\lambda t}\,dt
  =\frac2\lambda\Expect[T]
  =\frac2{\lambda^2}.
\]
We will use this second moment to measure spread.

\subsection{Why expectations add}

Suppose completing a job requires a wait \(X\) and a processing time \(Y\).
On each repetition, the total is \(X+Y\). Averaging the totals gives the
average waiting time plus the average processing time. This remains true
even if long waits tend to accompany long processing times.

\begin{theorem}[Linearity of expectation]
Let \(X_1,\ldots,X_n\) have finite expectations.
For constants \(a_i\) and \(b\),
\[
  \Expect\!\left[b+\sum_{i=1}^n a_iX_i\right]
  =b+\sum_{i=1}^n a_i\Expect[X_i].
\]
Independence is not required.
\end{theorem}

For a finite sample space, distributing the sum makes the reason explicit:
\[
  \sum_{\omega\in\Omega}
    (X(\omega)+Y(\omega))\Prob(\{\omega\})
  =\Expect[X]+\Expect[Y].
\]
The integral version works in the same way.

For instance, if a job has a fixed two-minute setup followed by
\(T\sim\operatorname{Exp}(0.2)\), its mean duration is
\(2+\Expect[T]=7\) minutes. Expressing \(T\) in seconds instead multiplies
its mean by \(60\).

Linearity also makes counts easy to analyze. For an event \(A\),
the indicator \(\mathbf1_A\) is one if \(A\) occurs and zero otherwise.
Its expectation is
\[
  \Expect[\mathbf1_A]=1\cdot\Prob(A)+0\cdot\Prob(A^c)=\Prob(A).
\]
Consequently, the expected number of events occurring among
\(A_1,\ldots,A_n\) is
\[
  \Expect\!\left[\sum_{i=1}^n\mathbf1_{A_i}\right]
  =\sum_{i=1}^n\Prob(A_i).
\]
No independence assumption is needed for this mean. Dependence will matter
when we ask about the count's spread.

\section{Variance and standard deviation}
\label{sec:variance}

A constant five-minute wait and an exponential wait with mean five minutes
have the same expectation. Their behavior is very different: the first
is perfectly predictable, while the second sometimes ends almost
immediately and sometimes takes much longer. We need a measure of this
variation around the mean.

If \(\mu=\Expect[X]\), the deviation of an observation is \(X-\mu\).
Averaging the signed deviations gives zero:
\(\Expect[X-\mu]=0\). Positive and negative deviations cancel.
Squaring them prevents that cancellation and gives more weight to large
departures from the mean.

\begin{definition}[Variance and standard deviation]
For \(X\) with a finite second moment,
\[
  \Var(X)=\Expect[(X-\mu)^2],\qquad \mu=\Expect[X],
  \qquad
  \operatorname{sd}(X)=\sqrt{\Var(X)}.
\]
\end{definition}

Variance is the average squared deviation. If waiting time is in minutes,
variance is in minutes squared. Taking the square root returns to minutes;
this is why standard deviation is useful when describing the size of
typical fluctuations. It is not a bound on individual observations.

Expanding the square gives a convenient way to calculate variance:
\begin{align}
  \Var(X)
  &=\Expect[X^2-2\mu X+\mu^2]\notag\\
  &=\Expect[X^2]-2\mu\Expect[X]+\mu^2
   =\Expect[X^2]-\Expect[X]^2.
  \label{eq:variance-moments}
\end{align}
For exponential waiting times, we have already computed both terms:
\[
  \Var(T)=\frac2{\lambda^2}-\frac1{\lambda^2}
  =\frac1{\lambda^2},
  \qquad \operatorname{sd}(T)=\frac1\lambda.
\]
The inbox with rate \(0.2\) has both mean and standard deviation equal to
five minutes. A constant five-minute wait instead has variance zero.

\subsection{Changing the origin or the units}

Adding a fixed delay \(b\) moves every observation and its mean by the
same amount. The deviations from the mean stay unchanged.
Multiplying times by \(a\) multiplies every deviation by \(a\), so their
squares are multiplied by \(a^2\). Therefore
\[
  \Expect[aX+b]=a\Expect[X]+b,\qquad
  \Var(aX+b)=a^2\Var(X).
\]
For example, converting minutes to seconds multiplies the mean and
standard deviation by \(60\), and the variance by \(60^2\).
Variance is not linear.

\subsection{Adding two variable delays}

For a job with waiting time \(X\) and processing time \(Y\), let
\(U=X-\Expect[X]\) and \(V=Y-\Expect[Y]\). The deviation of the total
duration is \(U+V\). Squaring gives
\[
  (U+V)^2=U^2+V^2+2UV.
\]
The first two terms describe the separate spreads. The last term tells us
whether the two delays fluctuate together. If both are above their means,
or both are below, their product \(UV\) is positive. If one is above its
mean while the other is below, the product is negative.

The average of this product is called the \emph{covariance}:
\[
  \Cov(X,Y)=\Expect[(X-\Expect[X])(Y-\Expect[Y])].
\]
Taking expectations in the squared-sum identity gives
\begin{equation}
  \Var(X+Y)=\Var(X)+\Var(Y)+2\Cov(X,Y).
  \label{eq:variance-sum}
\end{equation}
For independent variables, expectations of products factor. Thus
\(\Expect[UV]=\Expect[U]\Expect[V]=0\), and the covariance term vanishes.
For independent variables with finite second moments, this extends to
\[
  \Var\!\left(\sum_{i=1}^n a_iX_i\right)
  =\sum_{i=1}^n a_i^2\Var(X_i).
\]

\begin{example}[Three different ways to combine fluctuations]
Let \(X,Y\) each have variance one. If they are independent, the variance
of \(X+Y\) is two. If \(Y=X\), the fluctuations reinforce each other and
\(\Var(X+Y)=\Var(2X)=4\). If \(Y=-X\), they cancel exactly and the sum
has variance zero.
\end{example}

Zero covariance is enough for variances to add, even without full
independence. We will study that distinction further when considering
joint distributions. For the sums used below, independence will be an
explicit part of the model.

\section{Poisson counts}
\label{sec:poisson}

Return to the inbox, but change the question. Instead of timing the next
message, watch for ten minutes and count how many messages arrive.
A waiting time can be \(2.4\) minutes; a message count must be an integer
such as zero, one, or two. The same arrivals produce two different kinds
of random variables.

Write \(N(t)\) for the number of messages arriving between time zero and
time \(t\). Figure~\ref{fig:arrival-timeline} shows one possible record.
At minute six we have counted two messages; at minute ten we have counted
three.

\begin{figure}[htbp]
  \centering
  \begin{tikzpicture}[x=1.15cm,y=1cm,font=\small]
  \fill[CourseBlue!7] (0,-0.16) rectangle (6,0.16);
  \draw[->] (0,0) -- (10.35,0);
  \foreach \t in {0,2,4,6,8,10}{
    \draw (\t,0.07) -- (\t,-0.07);
    \node[below] at (\t,-0.11) {\(\t\)};
  }
  \draw[CourseGray,dashed] (6,-0.12) -- (6,1.0);
  \node[above,CourseBlue] at (6,1.0) {\(t=6,\quad N(6)=2\)};
  \foreach \t/\k in {1.5/1,4/2,8/3}{
    \fill[CourseBlue] (\t,0) circle (2.5pt);
    \node[above] at (\t,0.2) {\(T_{\k}\)};
  }
  \draw[<->,CourseBlue] (0,-0.85) -- (1.5,-0.85);
  \node[below] at (0.75,-0.85) {\(W_1=1.5\)};
  \draw[<->,CourseBlue] (1.5,-0.85) -- (4,-0.85);
  \node[below] at (2.75,-0.85) {\(W_2=2.5\)};
  \draw[<->,CourseBlue] (4,-0.85) -- (8,-0.85);
  \node[below] at (6,-0.85) {\(W_3=4\)};
  \node[below,align=center] at (9.35,-0.85) {Time\\(minutes)};
\end{tikzpicture}

  \caption{One possible sequence of arrivals, at \(1.5\), \(4\), and \(8\)
  minutes. The gaps \(W_1,W_2,W_3\) are waiting times; the arrival times
  \(T_k\) are cumulative sums of those gaps. By the marked time \(t=6\),
  two messages have arrived, so \(N(6)=2\).}
  \label{fig:arrival-timeline}
\end{figure}

\subsection{Counting successful short intervals}

We can use the same short-interval model as before. Divide \(t\) minutes
into \(n\) tiny intervals, each independently containing a message with
probability \(p_n=\lambda t/n\). To get exactly \(k\) messages, exactly
\(k\) intervals must contain a message and the others must be empty.

For a particular choice of those \(k\) intervals, independence gives
probability \(p_n^k(1-p_n)^{n-k}\). The number of ways to choose the
intervals is
\[
  \binom nk=\frac{n!}{k!(n-k)!}.
\]
For example, if there are three intervals and exactly one message, it can
be in the first, second, or third interval: three possibilities.
Here \(n!=1\cdot2\cdots n\), with \(0!=1\).
Adding the equally probable possibilities gives
\[
  \binom nk\left(\frac{\lambda t}{n}\right)^k
       \left(1-\frac{\lambda t}{n}\right)^{n-k}.
\]

Put \(m=\lambda t\). For fixed \(k\), as \(n\) increases,
\[
  \frac{\binom nk}{n^k}\longrightarrow\frac1{k!},
  \qquad
  \left(1-\frac mn\right)^{n-k}\longrightarrow e^{-m}.
\]
Thus the continuous-time count has probability
\(e^{-m}m^k/k!\) of equalling \(k\). This is the \emph{Poisson distribution}.

\begin{definition}[Poisson distribution]
For \(m>0\), a count \(N\sim\operatorname{Poisson}(m)\) has PMF
\[
  \Prob(N=k)=e^{-m}\frac{m^k}{k!},
  \qquad k=0,1,2,\ldots.
\]
For our inbox observed for \(t\) minutes,
\(N(t)\sim\operatorname{Poisson}(\lambda t)\).
\end{definition}

The rate \(\lambda\) describes arrivals per minute; \(m=\lambda t\)
describes the expected number over the whole observation window.
At rate \(0.2\) per minute, a ten-minute window has \(m=2\), while a
thirty-minute window has \(m=6\).
Neither count is guaranteed to equal its mean.

\begin{example}[No messages or several messages]
For the ten-minute count with \(m=2\),
\[
  \Prob(N(10)=0)=e^{-2}\approx0.1353,
\]
and
\[
  \Prob(N(10)\ge2)
  =1-\Prob(N(10)=0)-\Prob(N(10)=1)
  =1-e^{-2}(1+2)\approx0.5940.
\]
No messages in ten minutes is exactly the event that the first waiting
time exceeds ten minutes. The Poisson and Exponential calculations agree:
\(\Prob(N(10)=0)=\Prob(T>10)=e^{-2}\).
\end{example}

\subsection{Mean and variance of the count}

To check normalization, sum the PMF and use the exponential series:
\[
  \sum_{k=0}^\infty e^{-m}\frac{m^k}{k!}=e^{-m}e^m=1.
\]
For the mean, the factor \(k\) cancels one factor in \(k!\).
Reindexing the remaining sum with \(j=k-1\) gives
\[
  \Expect[N]
  =\sum_{k=1}^\infty k e^{-m}\frac{m^k}{k!}
  =m\sum_{j=0}^\infty e^{-m}\frac{m^j}{j!}=m.
\]

For the variance we need \(\Expect[N^2]\). It is easier to first calculate
\(\Expect[N(N-1)]\), since \(k(k-1)\) cancels two factors in \(k!\):
\[
  \Expect[N(N-1)]
  =m^2\sum_{j=0}^\infty e^{-m}\frac{m^j}{j!}=m^2.
\]
Now use \(N^2=N(N-1)+N\):
\[
  \Expect[N^2]=m^2+m,\qquad \Var(N)=m.
\]
For ten minutes, the mean and variance are both two, and the standard
deviation is \(\sqrt2\) messages.

The equality of mean and variance will later help us recognize when a
single Poisson model misses part of the variation in a population.

\subsection{Combining independent arrival streams}

Suppose two independent inboxes receive messages at rates \(\lambda_1\)
and \(\lambda_2\). Counting all their messages over the same interval
should add the expected counts. In fact, the entire distribution remains
Poisson:
\[
  N_1+N_2\sim\operatorname{Poisson}(m_1+m_2),
  \qquad m_i=\lambda_i t.
\]
To see why, split a total of \(k\) messages into \(j\) from the first inbox
and \(k-j\) from the second. Independence lets us multiply their
probabilities; summing over \(j\) and using the binomial theorem gives
\[
  \Prob(N_1+N_2=k)
  =\frac{e^{-(m_1+m_2)}}{k!}
    \sum_{j=0}^k\binom kj m_1^j m_2^{k-j}
  =\frac{e^{-(m_1+m_2)}(m_1+m_2)^k}{k!}.
\]
In the short-interval model, counts in disjoint time windows are also
independent. If arrivals instead trigger further arrivals or the rate
changes with the time of day, these assumptions need to be reconsidered.

\section{Gamma waiting times and the event-count connection}
\label{sec:gamma}

Suppose the help desk collects three messages before starting a batch.
Now we need the time until the third message, rather than the time until
the next one or the number received by a fixed deadline.

Let \(W_1\) be the wait from time zero to the first message, \(W_2\) the
gap from the first to the second, and so on. In our arrival model these
gaps are independent and each follows \(\operatorname{Exp}(\lambda)\).
The \(k\)-th arrival time is
\[
  T_k=W_1+\cdots+W_k.
\]
For example, the gaps in Figure~\ref{fig:arrival-timeline} are \(1.5\),
\(2.5\), and \(4\) minutes, so \(T_3=8\) minutes.

Before deriving a new density, we can already find the mean and variance:
\[
  \Expect[T_k]=\frac{k}{\lambda},\qquad
  \Var(T_k)=\frac{k}{\lambda^2}.
\]
Expectation adds by linearity; variance adds because the gaps are
independent. At rate \(0.2\), waiting for three messages takes \(15\)
minutes on average, with variance \(75\) minutes squared.

\subsection{The distribution of a sum of waiting times}

Consider two messages first. A total wait near five minutes could consist
of a first gap near one minute and a second near four, or a first near
two and a second near three, and so on. To find the density of the total,
we must account for every possible split.

If the first gap is \(u\), the second must be near \(t-u\).
Independence gives a product of densities for these two contributions.
Integrating over \(0<u<t\) adds all the possible splits:
\begin{align*}
  f_{T_2}(t)
  &=\int_0^t f_{W_1}(u)f_{W_2}(t-u)\,du\\
  &=\int_0^t\lambda e^{-\lambda u}
       \lambda e^{-\lambda(t-u)}\,du
   =\lambda^2 t e^{-\lambda t},\qquad t>0.
\end{align*}
This way of finding the density of a sum is called \emph{convolution}.

Repeating the same step for \(k\) independent gaps gives
\[
  f_{T_k}(t)=\frac{\lambda^k}{(k-1)!}t^{k-1}e^{-\lambda t},
  \qquad t>0.
\]
For \(k=1\), this is the exponential density. For \(k>1\), the density
starts near zero: fitting several arrivals into a very short time is
unlikely. As \(t\) increases, enough time becomes available for all
\(k\) arrivals. Eventually the density decreases because exceptionally
long total waits are also unlikely.

\begin{figure}[htbp]
  \centering
  \begin{tikzpicture}
\begin{axis}[
  width=0.94\textwidth,height=6.1cm,
  axis lines=left,xmin=0,xmax=10,ymin=0,ymax=1.08,
  xlabel={Waiting time \(t\)},ylabel={Density},
  tick label style={font=\small},label style={font=\small},
  legend style={draw=none,font=\small,at={(0.98,0.98)},anchor=north east},
  grid=major,grid style={black!8},domain=0:10,samples=160,
]
  \addplot[CourseBlue,very thick] {exp(-x)};
  \addlegendentry{\(\operatorname{Exp}(1)=\operatorname{Gamma}(1,1)\)}
  \addplot[orange!80!black,thick,dashed] {x*exp(-x)};
  \addlegendentry{\(\operatorname{Gamma}(2,1)\)}
  \addplot[teal!80!black,thick,dashdotted] {x^3*exp(-x)/6};
  \addlegendentry{\(\operatorname{Gamma}(4,1)\)}
\end{axis}
\end{tikzpicture}

  \caption{Time until the first, second, or fourth arrival, at rate one.
  One arrival can occur almost immediately. Several arrivals require a sum
  of waiting times, moving the typical total wait to the right.
  The means are \(1,2,4\); the standard deviations are \(1,\sqrt2,2\).}
  \label{fig:waiting-times}
\end{figure}

\subsection{From an integer number of arrivals to a shape parameter}

The density above gives a useful family of positive, often skewed
distributions. We can extend its shape beyond integer values of \(k\).
The only obstacle is the factorial \((k-1)!\). Its continuous extension
is the gamma function:
\[
  \Gamma(\alpha)=\int_0^\infty u^{\alpha-1}e^{-u}\,du,\qquad\alpha>0.
\]
Integration by parts gives
\(\Gamma(\alpha+1)=\alpha\Gamma(\alpha)\). Since \(\Gamma(1)=1\), this
implies \(\Gamma(k)=(k-1)!\) for positive integers \(k\).

Replacing \(k\) by a positive real number \(\alpha\) and the factorial by
\(\Gamma(\alpha)\) gives the \emph{Gamma distribution}. We write its rate
as \(\beta\) so that \(\lambda\) can continue to denote the particular
inbox rate.

\begin{definition}[Gamma distribution]
For shape \(\alpha>0\) and rate \(\beta>0\),
\(X\sim\operatorname{Gamma}(\alpha,\beta)\) has density
\[
  f_X(x)=\frac{\beta^\alpha}{\Gamma(\alpha)}
      x^{\alpha-1}e^{-\beta x},\qquad x>0,
\]
and zero density for \(x<0\).
\end{definition}

The substitution \(u=\beta x\) shows that its integral is one.
The case \(\alpha=1\) is Exponential; integer \(\alpha=k\) describes
waiting for \(k\) arrivals at rate \(\beta\).
For noninteger shape, the density still makes sense, but there is no
literal fractional-number-of-arrivals interpretation.

We can calculate the moments for every positive shape by inserting
\(x\) or \(x^2\) into the density integral:
\begin{align*}
  \Expect[X]
    &=\frac{\Gamma(\alpha+1)}{\beta\Gamma(\alpha)}
      =\frac{\alpha}{\beta},\\
  \Expect[X^2]
    &=\frac{\Gamma(\alpha+2)}{\beta^2\Gamma(\alpha)}
      =\frac{\alpha(\alpha+1)}{\beta^2}.
\end{align*}
Subtracting the squared mean gives
\[
  \Var(X)=\frac{\alpha}{\beta^2}.
\]
Increasing the rate \(\beta\) shortens the waits.
Increasing the shape \(\alpha\) at fixed rate increases the mean, while
the spread relative to that mean decreases:
\(\operatorname{sd}(X)/\Expect[X]=1/\sqrt{\alpha}\).

Throughout these notes the second Gamma parameter is a rate.
Another common convention uses a scale \(\theta=1/\beta\), measured in
the same units as \(X\). In that convention the mean and variance are
\(\alpha\theta\) and \(\alpha\theta^2\). Checking which parameterization
is used prevents a reciprocal error.

\subsection{A deadline question can be answered by counting}

To have received the third message by minute ten means to have counted
at least three messages by minute ten. These are the same event,
expressed using different random variables:
\begin{equation}
  \{T_k\le t\}=\{N(t)\ge k\},\qquad k\ge1.
  \label{eq:arrivals-counts}
\end{equation}
Therefore, instead of integrating a Gamma density, we can add Poisson
probabilities:
\[
  \Prob(T_k\le t)
  =1-\Prob(N(t)<k)
  =1-e^{-\lambda t}\sum_{j=0}^{k-1}\frac{(\lambda t)^j}{j!}.
\]
For rate \(0.2\), a ten-minute window has expected count two, so
\[
  \Prob(T_3\le10)
  =1-e^{-2}(1+2+2^2/2)\approx0.3233.
\]
Receiving exactly \(k\) messages also has a waiting-time description:
the \(k\)-th has arrived, but the \((k+1)\)-th has not.
For \(k\ge1\),
\(\Prob(N(t)=k)=\Prob(T_k\le t)-\Prob(T_{k+1}\le t)\).
For \(k=0\), it is simply \(\Prob(T_1>t)\).

\section{Bernoulli and Binomial distributions}
\label{sec:bernoulli-binomial}

The inbox model asked whether a message arrived during each small time
interval. Many other observations are also binary: a link was clicked
or it was not, a transmission succeeded or failed, an item passed an
inspection or failed it. We now describe these trials without reference
to time.

\subsection{One binary outcome}

Code a click as one and no click as zero. If the click probability is
\(p\), then
\[
  \Prob(B=1)=p,\qquad \Prob(B=0)=1-p,\qquad 0\le p\le1.
\]
This is a \emph{Bernoulli distribution}, written
\(B\sim\operatorname{Bernoulli}(p)\).
The same construction applies to the indicator of any event.

For a binary variable, squaring changes neither possible value:
\(0^2=0\), \(1^2=1\). Thus
\[
  \Expect[B]=p,\qquad \Expect[B^2]=p,\qquad \Var(B)=p(1-p).
\]
At \(p=1/2\), neither outcome is favored and the variance is largest.
At \(p=0\) or \(p=1\), every observation is the same and variance is zero.

\subsection{A fixed number of independent trials}

Show a link to \(n\) people. Assume their clicks are independent and each
has the same probability \(p\). If \(B_i\) indicates the \(i\)-th click,
then the total number of clicks is
\[
  S_n=B_1+\cdots+B_n.
\]
It follows a \emph{Binomial distribution},
\(S_n\sim\operatorname{Binomial}(n,p)\), with
\[
  \Prob(S_n=k)=\binom nk p^k(1-p)^{n-k},\qquad k=0,\ldots,n.
\]
The counting argument is the one used for the short arrival intervals:
choose the \(k\) successful trials, multiply their success and failure
probabilities, and add over all choices. The displayed expression is
read directly for \(0<p<1\); at \(p=0\) or \(p=1\) the count is fixed.

We do not need to sum this PMF to obtain its moments. Since the count is
a sum of independent Bernoulli variables,
\[
  \Expect[S_n]=np,\qquad \Var(S_n)=np(1-p).
\]
For \(n\ge1\), the observed click proportion is \(\widehat p=S_n/n\).
The scaling rules for mean and variance give
\[
  \Expect[\widehat p]=p,\qquad
  \Var(\widehat p)=\frac{p(1-p)}{n}.
\]
With more independent observations, the proportion fluctuates less.
For example, multiplying \(n\) by four halves its standard deviation.

\begin{example}[A small click probability]
For \(100\) independent impressions with \(p=0.03\), the count has mean
\(3\) and variance \(2.91\). The chance of no clicks is
\[
  \Prob(S_{100}=0)=(1-0.03)^{100}\approx0.0476.
\]
The common-\(p\) assumption treats all impressions alike. We will revisit
it when different users have different propensities to click.
\end{example}

\subsection{Why rare Binomial successes look Poisson}

The short-interval model contained many independent trials, each with a
tiny success probability. A Binomial count behaves similarly whenever
\(n\) is large and \(p\) is small, with \(np=m\) held fixed.
For any fixed \(k\),
\[
  \binom nk\left(\frac mn\right)^k
    \left(1-\frac mn\right)^{n-k}
  \longrightarrow e^{-m}\frac{m^k}{k!}.
\]
This is the Poisson limit derived earlier. Its moments agree with the
Binomial moments in the limit:
\[
  np=m,\qquad np(1-p)\longrightarrow m.
\]
For the click example, the approximation is
\(\operatorname{Poisson}(3)\), which gives no-click probability
\(e^{-3}\approx0.0498\), close to the exact \(0.0476\).

\section{Geometric waiting times}
\label{sec:geometric}

Now show the link repeatedly until the first click, instead of fixing
the number of impressions in advance. Let \(G\) be the number of
impressions needed, including the successful one.
For example, the sequence
\[
  \text{no click},\quad\text{no click},\quad\text{click}
\]
gives \(G=3\).

Assume independent trials with the same success probability \(0<p\le1\).
For \(0<p<1\), the event \(G=k\) requires \(k-1\) failures followed by
one success:
\[
  \Prob(G=k)=(1-p)^{k-1}p,\qquad k=1,2,\ldots.
\]
This is the \emph{Geometric distribution},
\(G\sim\operatorname{Geometric}(p)\).
If \(p=1\), the first trial always succeeds and \(G=1\).

\subsection{Waiting longer means accumulating failures}

The event \(G>k\) means that all of the first \(k\) trials failed. Hence
\[
  \Prob(G>k)=(1-p)^k,\qquad k=0,1,\ldots.
\]
With \(p=0.1\), the chance of still waiting for a success after twenty
trials is \(0.9^{20}\approx0.1216\).

This model has the same memoryless logic as exponential waiting.
Once the first \(m\) trials have failed, the next trials still have the
original success probability. Formally, for \(0<p<1\),
\[
  \Prob(G>m+k\mid G>m)
  =\frac{(1-p)^{m+k}}{(1-p)^m}
  =(1-p)^k.
\]
An example is a scrolling session in which a user independently stops
after each video with probability \(p\). The number of videos watched is
Geometric. If fatigue makes stopping more likely after each video,
the constant-\(p\) assumption no longer holds.

\subsection{Mean and variance}

To compute the mean, weight each possible stopping trial by its
probability. Write \(q=1-p\). The geometric series
\[
  \sum_{k=0}^{\infty}q^k=\frac1{1-q},\qquad |q|<1,
\]
can be differentiated to obtain the factor \(k\) needed in the mean:
\[
  \sum_{k=1}^{\infty}kq^{k-1}=\frac1{(1-q)^2}.
\]
Multiplying by \(p\) gives \(\Expect[G]=1/p\).
A second differentiation gives
\[
  \Expect[G(G-1)]
  =p\sum_{k=2}^{\infty}k(k-1)q^{k-1}
  =\frac{2(1-p)}{p^2}.
\]
As in the Poisson calculation, use \(G^2=G(G-1)+G\) to get
\[
  \Var(G)=\frac{1-p}{p^2}.
\]
At \(p=0.1\), the mean is ten trials and the variance is ninety.

Some sources count failures before success instead. That variable is
\(H=G-1\), so its mean is \((1-p)/p\), its variance is unchanged,
and it can equal zero. In these notes, Geometric always counts the
successful trial too.

\subsection{Turning a trial count into a time}

Return to the inbox approximation, where one trial corresponds to a
short interval of duration \(\Delta\) minutes.
A count of \(G\) intervals represents a waiting time of \(\Delta G\)
minutes. Set \(p=\lambda\Delta\), just as in the first section.

By time \(t\), we have completed \(\lfloor t/\Delta\rfloor\) intervals;
the floor symbol means the integer number of whole intervals.
No arrival by that time has probability
\[
  \Prob(\Delta G>t)
  =(1-\lambda\Delta)^{\lfloor t/\Delta\rfloor}
  \longrightarrow e^{-\lambda t}
  \qquad(\Delta\downarrow0).
\]
The Geometric waiting count therefore becomes an Exponential waiting
time as the time intervals shrink.
The moments show the same connection:
\[
  \Expect[\Delta G]=\frac1\lambda,\qquad
  \Var(\Delta G)=\frac{1-\lambda\Delta}{\lambda^2}
  \longrightarrow\frac1{\lambda^2}.
\]
The time-unit conversion is what connects the two distributions:
Geometric counts trials; Exponential measures elapsed time.

\section{The Gaussian distribution}
\label{sec:gaussian}

The waiting times studied so far are nonnegative and usually skewed:
a wait can be much longer than average, but cannot extend below zero.
Measurement errors pose a different problem. A sensor measuring a fixed
quantity may sometimes read too high and sometimes too low. Several
small disturbances can add together, with positive and negative
contributions partly cancelling.

A useful model for such errors has a central peak and decreases
symmetrically on either side. This is the \emph{Gaussian}, also called
the \emph{normal}, distribution. It is often a useful approximation when
many small independent disturbances contribute and no one disturbance
dominates. The connection between sums and Gaussian approximations will
be developed further in the chapter on limit theorems.

\subsection{Center and width}

Let \(\mu\) describe the center of the errors and \(\sigma>0\) their
scale. The center can be nonzero if the sensor is systematically biased.
The density is
\[
  f_X(x)=\frac1{\sigma\sqrt{2\pi}}
    \exp\!\left[-\frac12\left(\frac{x-\mu}{\sigma}\right)^2\right],
  \qquad x\in\Real.
\]
We write \(X\sim\mathcal N(\mu,\sigma^2)\).

The expression \((x-\mu)/\sigma\) measures how far \(x\) lies from the
center relative to the width of the curve. Squaring it treats positive
and negative deviations equally. The negative exponential makes values
far from the center less likely. Increasing \(\sigma\) widens the curve;
the factor \(1/\sigma\) reduces its height to preserve total area one.

The simplest member is \(Z\sim\mathcal N(0,1)\), with density
\[
  \phi(z)=\frac1{\sqrt{2\pi}}e^{-z^2/2}.
\]
The constant \(\sqrt{2\pi}\) is the value of
\(\int_{-\infty}^{\infty}e^{-z^2/2}\,dz\), so this density integrates to
one. Every other Gaussian can be obtained by shifting and scaling:
\[
  X=\mu+\sigma Z.
\]

\subsection{Why the parameters are the mean and standard deviation}

Symmetry of \(\phi\) means that the contributions of \(z\) and \(-z\)
to the mean cancel, giving \(\Expect[Z]=0\).
For the second moment, differentiating the density gives
\(\phi'(z)=-z\phi(z)\). This turns the moment integral into an
integration-by-parts calculation:
\begin{align*}
  \Expect[Z^2]
  &=\int_{-\infty}^{\infty}z^2\phi(z)\,dz
    =-\int_{-\infty}^{\infty}z\phi'(z)\,dz\\
  &=[-z\phi(z)]_{-\infty}^{\infty}
      +\int_{-\infty}^{\infty}\phi(z)\,dz
    =1.
\end{align*}
Thus \(\Var(Z)=1\). Applying the shift and scale rules to
\(X=\mu+\sigma Z\) gives
\[
  \Expect[X]=\mu,\qquad \Var(X)=\sigma^2,\qquad
  \operatorname{sd}(X)=\sigma.
\]
This justifies the parameters' interpretation.
In \(\mathcal N(\mu,\sigma^2)\), the second argument is the variance.

\subsection{Expressing an error in standard-deviation units}

For a sensor with mean error zero, an error of \(0.4\) units may be large
or small relative to its usual variation. If the standard deviation is
\(0.2\), that error is
two standard deviations above the mean; if it is \(2\), it is only
\(0.2\) standard deviations above the mean.

The transformation
\[
  Z=\frac{X-\mu}{\sigma}
\]
subtracts the center and expresses the deviation in standard-deviation
units. This is \emph{standardization}. For a Gaussian \(X\), the result
always follows \(\mathcal N(0,1)\). We can therefore answer interval
questions for any Gaussian using a single CDF:
\[
  \Phi(z)=\Prob(Z\le z)=\int_{-\infty}^z\phi(u)\,du.
\]
The endpoints of the original interval are transformed in the same way:
\[
  \Prob(a<X\le b)
  =\Phi\!\left(\frac{b-\mu}{\sigma}\right)
   -\Phi\!\left(\frac{a-\mu}{\sigma}\right).
\]
The integral defining \(\Phi\) is evaluated numerically. Symmetry gives
\(\Phi(-z)=1-\Phi(z)\).

\begin{example}[How often does a sensor make a large error?]
Suppose the error is \(E\sim\mathcal N(0,0.2^2)\). An absolute error
larger than \(0.4\) is an error more than two standard deviations from
zero. The two symmetric tails give
\[
  \Prob(|E|>0.4)
  =\Prob(|Z|>2)
  =2[1-\Phi(2)]\approx0.0455.
\]
Under this model, about \(4.55\%\) of repeated readings have such an error.
Likewise,
\[
  \Prob(|Z|\le1)\approx0.6827,\qquad
  \Prob(|Z|\le1.96)\approx0.95.
\]
The central \(95\%\) interval for the sensor error is therefore
\[
  [-1.96(0.2),1.96(0.2)]=[-0.392,0.392].
\]
\end{example}

\subsection{Adding errors and connecting back to earlier distributions}

Changing units or adding a fixed bias preserves the Gaussian shape:
\[
  aX+b\sim\mathcal N(a\mu+b,a^2\sigma^2),\qquad a\ne0.
\]
For \(a=0\), the result is simply the constant \(b\).

Independent Gaussian errors also add to another Gaussian. If
\(X\sim\mathcal N(\mu_X,\sigma_X^2)\) and
\(Y\sim\mathcal N(\mu_Y,\sigma_Y^2)\) are independent, then
\[
  X+Y\sim\mathcal N(\mu_X+\mu_Y,\sigma_X^2+\sigma_Y^2).
\]
The convolution method used for Gamma waiting times gives this density.
The means and variances agree with the addition rules we already know.

For example, two independent centered errors with standard deviations
\(0.2\) and \(0.3\) produce a total error with variance
\(0.2^2+0.3^2=0.13\), and standard deviation \(\sqrt{0.13}\approx0.361\).
Standard deviations do not simply add.

A sum can also be approximately Gaussian even when its individual terms
are not Gaussian. We have already met two such sums: Binomial counts
add Bernoulli outcomes, and integer-shape Gamma waits add exponential
gaps. For sufficiently large parameters, useful approximations are
\[
  \begin{array}{rcl}
    \operatorname{Binomial}(n,p)
      &\approx&\mathcal N(np,np(1-p)),\\
    \operatorname{Gamma}(k,\lambda)
      &\approx&\mathcal N(k/\lambda,k/\lambda^2),\\
    \operatorname{Poisson}(m)
      &\approx&\mathcal N(m,m).
  \end{array}
\]
Each Gaussian uses the mean and variance of the distribution being
approximated. For the Binomial, both the expected successes \(np\)
and failures \(n(1-p)\) should be large; for Gamma and Poisson, the
approximations improve as \(k\) and \(m\) grow.

The Gaussian allows negative values, while counts and waiting times do
not. Thus approximating a very small count or a strongly skewed waiting
time by a Gaussian can assign substantial probability to impossible
outcomes. The original distributions remain the appropriate starting
point in those cases.

\section{Mixtures and heterogeneity}
\label{sec:mixtures}

Suppose two sensors measure the same reference quantity. One tends to
read two units too low; the other tends to read two units too high.
Within either sensor, the random error has standard deviation one.

Choose a sensor with equal probability for each reading. If we combine
all readings without keeping track of which sensor produced them, what
distribution do the errors follow? Averaging the two sensor biases gives
zero, but that does not mean either sensor usually makes an error near
zero. Most errors are still near \(-2\) or \(2\).

\subsection{Choose a source, then draw an observation}

Let \(Z\) record which sensor is selected:
\[
  \Prob(Z=1)=\Prob(Z=2)=\frac12.
\]
Given the chosen sensor, the error distribution is
\[
  X\mid Z=1\sim\mathcal N(-2,1),\qquad
  X\mid Z=2\sim\mathcal N(2,1).
\]
The notation \(X\mid Z=k\) means that we are considering only readings
from sensor \(k\).

Write the two Gaussian densities as \(f_1\) and \(f_2\). For any interval,
half the readings come from each sensor, so the law of total probability
gives half of each sensor's interval probability. The density is therefore
\[
  f_X(x)=\frac12f_1(x)+\frac12f_2(x).
\]
This is a \emph{mixture distribution}. Its two peaks reflect the two
sources of readings.

The same construction works with \(K\) sources chosen with probabilities
\(\pi_1,\ldots,\pi_K\):
\begin{equation}
  f_X(x)=\sum_{k=1}^K\pi_k f_k(x),
  \qquad \pi_k\ge0,\quad \sum_{k=1}^K\pi_k=1.
  \label{eq:finite-mixture}
\end{equation}
The \(f_k\) are called the mixture's \emph{components}.
For discrete observations, use component PMFs in the same formula.
If the source label is not recorded, it is an unobserved, or
\emph{latent}, variable.

\subsection{Where the pooled variance comes from}

In the two-sensor example the overall mean is
\[
  \Expect[X]=\tfrac12(-2)+\tfrac12(2)=0.
\]
Within either sensor, the second moment equals its variance plus the
square of its mean:
\[
  \Expect[X^2\mid Z=1]=1+(-2)^2=5,\qquad
  \Expect[X^2\mid Z=2]=1+2^2=5.
\]
Averaging these gives \(\Expect[X^2]=5\). Hence the pooled variance is
\(5-0^2=5\), even though each sensor separately has variance one.

There are two sources of spread: random fluctuation within a sensor,
and the difference between the two sensors' centers. To separate them
in general, let component \(k\) have mean \(\mu_k\) and variance
\(\sigma_k^2\). Then
\[
  \mu=\Expect[X]=\sum_k\pi_k\mu_k,\qquad
  \Expect[X^2]=\sum_k\pi_k(\sigma_k^2+\mu_k^2).
\]
Subtracting \(\mu^2\) and regrouping the terms gives
\begin{equation}
  \boxed{
  \Var(X)
  =\underbrace{\sum_k\pi_k\sigma_k^2}_{\text{variation within sources}}
   +\underbrace{\sum_k\pi_k(\mu_k-\mu)^2}_{\text{variation between sources}}.
  }
  \label{eq:mixture-variance}
\end{equation}
For the sensors, these contributions are \(1\) and \(4\), respectively.

The inner averages can be written as conditional expectations.
The quantity \(\Expect[X\mid Z]\) equals \(\mu_k\) when source \(k\) is
selected; it varies because the selected source varies.
Similarly, \(\Var(X\mid Z)\) is the variance inside the selected source.
The outer expectation averages over the choice of source. With this
notation, the two identities become
\begin{align}
  \Expect[X]&=\Expect[\Expect[X\mid Z]],\label{eq:total-expectation}\\
  \Var(X)&=\Expect[\Var(X\mid Z)]
       +\Var(\Expect[X\mid Z]).\label{eq:total-variance}
\end{align}
These are the \emph{laws of total expectation and total variance}.
We have derived them for a finite mixture with finite second moments.

\subsection{The same mean and variance can hide different shapes}

A single Gaussian \(\mathcal N(0,5)\) has the same mean and variance as
the two-sensor mixture. Its most likely region is near zero, whereas
the mixture has two peaks near the sensor biases
(Figure~\ref{fig:gaussian-mixture}).

\begin{figure}[htbp]
  \centering
  \begin{tikzpicture}
\begin{axis}[
  width=0.94\textwidth,height=6.6cm,
  axis lines=left,xmin=-6,xmax=6,ymin=0,ymax=0.31,
  xlabel={Observation \(x\)},ylabel={Density},
  tick label style={font=\small},label style={font=\small},
  legend style={draw=none,font=\small,at={(0.5,0.98)},anchor=north},
  grid=major,grid style={black!8},domain=-6:6,samples=200,
]
  \addplot[CourseBlue,very thick]
    {0.5/sqrt(2*pi)*(exp(-(x+2)^2/2)+exp(-(x-2)^2/2))};
  \addlegendentry{Equal mixture of \(\mathcal N(-2,1)\) and \(\mathcal N(2,1)\)}
  \addplot[orange!80!black,thick,dashed] {exp(-x^2/10)/sqrt(10*pi)};
  \addlegendentry{Moment-matched \(\mathcal N(0,5)\)}
  \addplot[teal!80!black,thick,densely dotted,forget plot]
    {0.5*exp(-(x+2)^2/2)/sqrt(2*pi)};
  \addplot[teal!80!black,thick,densely dotted,forget plot]
    {0.5*exp(-(x-2)^2/2)/sqrt(2*pi)};
\end{axis}
\end{tikzpicture}

  \caption{Pooling two biased sensors produces the blue mixture density.
  The dotted curves are the weighted sensor densities, which add to the
  blue curve. A single Gaussian with the same mean and variance puts
  substantially more probability near zero.}
  \label{fig:gaussian-mixture}
\end{figure}

For example, in sensor 1 the event \(-1\le X\le1\) corresponds to a
standardized error between one and three. Sensor 2 gives the same
probability by symmetry. Therefore the pooled probability is
\[
  \Prob(|X|\le1)=\Phi(3)-\Phi(1)\approx0.1573.
\]
The Gaussian \(\mathcal N(0,5)\) instead gives
\(2\Phi(1/\sqrt5)-1\approx0.3453\).
Knowing the mean and variance alone does not determine interval
probabilities.

Also distinguish selecting one sensor from averaging two readings.
If we take one independent reading from \emph{each} sensor and average
them, their opposite biases cancel:
\[
  \frac{X_1+X_2}{2}\sim\mathcal N(0,1/2).
\]
A mixture chooses which source supplies an observation. An average
combines observations from both sources. These are different experiments.

\subsection{Different users have different message rates}

Two sensors give a small number of groups. In a user population, the
differences can vary continuously. Some users receive very few messages;
others receive many. This variation between members of a population is
called \emph{heterogeneity}.

Observe each user for one day. Let \(\Lambda\) be the user's expected
number of messages in that day. A particular user's value, denoted
\(\lambda\), is fixed during the observation; the uppercase variable
\(\Lambda\) varies when we choose a different user.
Assume that, for a user with value \(\lambda\),
\[
  N\mid\Lambda=\lambda\sim\operatorname{Poisson}(\lambda).
\]
Thus even for a fixed user the realized count \(N\) fluctuates from day
to day. Across users there is additional variation in the mean count.

The within-user mean and variance both equal \(\Lambda\). Averaging over
users, the laws of total expectation and variance give
\[
  \Expect[N]=\Expect[\Lambda],\qquad
  \Var(N)=\Expect[\Lambda]+\Var(\Lambda).
\]
The first contribution to the variance is daily Poisson variation.
The second is variation in the users' average activity.
If the latter is positive, the pooled variance exceeds the pooled mean.
This is called \emph{overdispersion} relative to a Poisson distribution.

Replacing every user's rate by the population average removes this
second contribution. It also misses the extra zeros contributed by
quiet users and the large counts contributed by busy users.

\subsection{A Gamma distribution of Poisson means}

We need a distribution for the positive quantity \(\Lambda\).
One possible choice is the Gamma family, whose shape and rate let us
describe different amounts of variation among users:
\[
  \Lambda\sim\operatorname{Gamma}(\alpha,\beta).
\]
The Gamma moments immediately give
\[
  \Expect[N]=\frac{\alpha}{\beta},\qquad
  \Var(N)=\frac{\alpha}{\beta}+\frac{\alpha}{\beta^2}.
\]
For \(\alpha=2,\beta=1/2\), the mean user count is four and
the variance between user means is eight. The pooled count therefore
has mean four and variance \(4+8=12\).

We can also find the full count distribution. Let \(g(\lambda)\) be the
Gamma density of user means. To find the probability of \(k\) messages,
average the Poisson probability over all possible user means:
\[
  \Prob(N=k)=\int_0^\infty
       \Prob(N=k\mid\Lambda=\lambda)\,g(\lambda)\,d\lambda.
\]
This is the continuous version of the weighted sum for the sensors:
possible user means replace the discrete source labels, and the integral
replaces the sum. The count remains discrete; it is the mixing parameter
that varies continuously.

Substituting the Poisson PMF and Gamma density gives
\begin{align*}
  \Prob(N=k)
  &=\int_0^\infty
     \frac{e^{-\lambda}\lambda^k}{k!}
     \frac{\beta^\alpha}{\Gamma(\alpha)}
       \lambda^{\alpha-1}e^{-\beta\lambda}\,d\lambda\\
  &=\frac{\Gamma(k+\alpha)}{\Gamma(\alpha)k!}
      \left(\frac{\beta}{\beta+1}\right)^\alpha
      \left(\frac1{\beta+1}\right)^k,\qquad k=0,1,\ldots.
\end{align*}
The second line follows by combining the powers and exponentials,
then using the gamma integral with shape \(k+\alpha\) and rate
\(\beta+1\). The resulting distribution is called \emph{negative binomial}.
Here its meaning is supplied by the construction: a Poisson count whose
mean differs from user to user according to a Gamma distribution.

For our numerical example,
\[
  \Prob(N=0)=\left(\frac{1/2}{1+1/2}\right)^2
            =\frac19\approx0.1111.
\]
A single \(\operatorname{Poisson}(4)\) distribution instead gives
\(e^{-4}\approx0.0183\). The mean counts agree, but the predicted fraction
of users receiving no messages is very different.

\begin{figure}[htbp]
  \centering
  \begin{tikzpicture}
\begin{axis}[
  width=0.94\textwidth,height=6.5cm,
  axis lines=left,xmin=-0.7,xmax=20.8,ymin=0,ymax=0.245,
  xlabel={Count \(k\)},ylabel={Probability mass},
  xtick={0,2,...,20},
  tick label style={font=\small},label style={font=\small},
  legend style={draw=none,font=\small,at={(0.98,0.98)},anchor=north east},
  grid=major,grid style={black!8},
]
  \addplot[CourseBlue,thick,ycomb,mark=*,mark size=1.9pt]
    coordinates {
      (0,0.01831564) (1,0.07326256) (2,0.14652511)
      (3,0.19536681) (4,0.19536681) (5,0.15629345)
      (6,0.10419563) (7,0.05954036) (8,0.02977018)
      (9,0.01323119) (10,0.00529248) (11,0.00192454)
      (12,0.00064151) (13,0.00019739) (14,0.00005640)
      (15,0.00001504) (16,0.00000376) (17,0.00000088)
      (18,0.00000020) (19,0.00000004) (20,0.00000001)
    };
  \addlegendentry{Poisson: mean \(4\), variance \(4\)}
  \addplot[orange!80!black,thick,ycomb,mark=square*,mark size=1.8pt,
    samples at={0,1,...,20},domain=0:20]
    {(x+1)/9*(2/3)^x};
  \addlegendentry{Gamma--Poisson: mean \(4\), variance \(12\)}
\end{axis}
\end{tikzpicture}

  \caption{Both populations average four messages per user per day.
  With a common Poisson mean of four, the variance is four.
  With Gamma-distributed user means, the variance is twelve, with more
  users receiving zero messages or unusually many messages.}
  \label{fig:heterogeneous-counts}
\end{figure}

The conditional-moment rules used here extend the finite-mixture
identities by replacing weighted sums with integrals. They require
finite moments, but they do not require us to evaluate the full mixture
PMF before calculating its mean and variance.

\subsection{Ten responses from one user or one response from ten users?}

Heterogeneity also matters for binary observations, but the way we group
observations changes the result. Suppose half the users click with
probability \(0.1\), and half with probability \(0.9\).
A randomly chosen user's response is a click with probability
\[
  \tfrac12(0.1)+\tfrac12(0.9)=0.5.
\]
That overall probability does not tell us how much the individual user
probabilities differ.

Now compare two experiments. In the first, choose one user and show
them ten posts, with independent responses once their click probability
is fixed. The chosen user is likely to click about once or about nine
times, depending on which group they belong to.
In the second, choose ten users independently and record one response
from each. Each response is independently a click with probability
\(0.5\), so the count is \(\operatorname{Binomial}(10,0.5)\).

The means are both five. The variances are different.
Write \(P\) for the first experiment's shared user probability and
\(S\) for the ten-response count. Within either user group, the variance
is \(10(0.1)(0.9)=0.9\). The group means are one and nine, whose variance
around five is sixteen. Therefore
\[
  \Var(S)=0.9+16=16.9,
\]
whereas the second experiment has variance \(10(0.5)(0.5)=2.5\).

More generally, let \(q=\Expect[P]\) and suppose
\(S\mid P\sim\operatorname{Binomial}(n,P)\). Then
\begin{align*}
  \Expect[S]&=nq,\\
  \Var(S)
    &=n\Expect[P(1-P)]+n^2\Var(P)\\
    &=nq(1-q)+n(n-1)\Var(P).
\end{align*}
The final step uses
\(\Expect[P(1-P)]=q(1-q)-\Var(P)\), obtained by expanding
\(\Expect[P^2]=\Var(P)+q^2\).
The extra variance comes from sharing the same random user probability
across trials.

If a fresh \(P_i\) is independently drawn for each trial instead, and the
responses are independently generated given those probabilities, the
responses are independent Bernoulli variables with parameter \(q\).
Their sum is \(\operatorname{Binomial}(n,q)\).
Which quantities are shared across observations is part of specifying
the experiment.

\section{Extension: Gaussian scale mixtures and Student's \texorpdfstring{\(t\)}{t}}
\label{sec:student}

Return to measurement errors. So far the two sensors differed in their
average errors. Another possibility is that a sensor is usually precise
but occasionally operates under much noisier conditions, without changing
its average error.

For example, suppose \(95\%\) of readings have error
\(\mathcal N(\mu,\sigma^2)\), and the remaining \(5\%\) have standard
deviation \(5\sigma\):
\[
  X\sim0.95\,\mathcal N(\mu,\sigma^2)
      +0.05\,\mathcal N(\mu,25\sigma^2).
\]
This is called a \emph{contaminated Gaussian} model.
The noisier readings make large errors more common. The mean stays
\(\mu\), while the mixture variance is
\[
  \Var(X)=0.95\sigma^2+0.05(25\sigma^2)=2.2\sigma^2.
\]

Instead of two noise levels, we can allow a continuous range.
Describe the noise using its inverse variance \(\tau\), often called
\emph{precision}. A small \(\tau\) means a large variance and a noisy
reading; a large \(\tau\) means a narrow error distribution.
For a standardized error, use
\[
  Z\mid\tau\sim\mathcal N(0,1/\tau).
\]

The Gamma family is a natural candidate for a positive random precision.
Choose
\[
  \tau\sim\operatorname{Gamma}(\nu/2,\nu/2),\qquad \nu>0.
\]
Using equal shape and rate fixes \(\Expect[\tau]=1\);
the Gamma variance formula gives \(\Var(\tau)=2/\nu\).
Thus \(\nu\) controls how much the noise level varies: large \(\nu\)
keeps precision close to one, while small \(\nu\) permits more variation.

Average the conditional Gaussian density over these precision values,
just as we averaged Poisson probabilities over user means.
Let \(g_\tau\) denote the Gamma density and \(f_{Z\mid\tau}(z\mid u)\)
the Gaussian density with variance \(1/u\). Then
\begin{align*}
  f_Z(z)&=\int_0^\infty f_{Z\mid\tau}(z\mid u)\,g_\tau(u)\,du\\
  &=\frac{\Gamma((\nu+1)/2)}
       {\sqrt{\nu\pi}\,\Gamma(\nu/2)}
    \left(1+\frac{z^2}{\nu}\right)^{-(\nu+1)/2}.
\end{align*}
The product inside the integral is proportional, as a function of \(u\),
to \(u^{(\nu-1)/2}e^{-(\nu+z^2)u/2}\).
This is another gamma integral, now with shape \((\nu+1)/2\)
and rate \((\nu+z^2)/2\).
The resulting density defines \emph{Student's \(t\)-distribution},
written \(Z\sim t_\nu\).
In this construction, \(\nu\) controls variation in precision and
therefore how often extreme errors occur.

Occasionally drawing a small precision produces an unusually large
Gaussian error. After averaging over precisions, extreme errors remain
much more likely than under a standard Gaussian. This is the meaning
of \emph{heavy tails} here: for large \(|z|\), the density decreases
like \(|z|^{-(\nu+1)}\), rather than like \(e^{-z^2/2}\).

For \(\nu>1\), the mean exists and is zero. For \(\nu>2\), total variance
and the Gamma integral give
\[
  \Var(Z)=\Expect[1/\tau]=\frac{\nu}{\nu-2}.
\]
When \(1<\nu\le2\), rare large observations make the second-moment
integral diverge, so the variance is infinite.
When \(0<\nu\le1\), even the absolute first moment diverges and the
mean is undefined. This illustrates why the convergence conditions in
the definitions of expectation and variance matter.

For an error with center \(\mu\) and scale \(s>0\), write \(X=\mu+sZ\).
For \(\nu>2\),
\[
  \Expect[X]=\mu,\qquad
  \Var(X)=s^2\frac{\nu}{\nu-2}.
\]
The scale \(s\) is therefore not its standard deviation.
As \(\nu\) increases, precision concentrates near one and the standard
Student distribution approaches \(\mathcal N(0,1)\).

\clearpage
\section{Distribution reference}
\label{sec:distribution-reference}

The table collects the means and variances derived above. The support lists
the values the variable can take, distinguishing counts from waiting times
and signed errors.

\begin{center}
\renewcommand{\arraystretch}{1.5}
\begin{tabularx}{\textwidth}{@{}L{0.31\textwidth}Ycc@{}}
\toprule
\textbf{Distribution} & \textbf{Support} &
  \(\boldsymbol{\Expect[X]}\) & \(\boldsymbol{\Var(X)}\)\\
\midrule
\(\operatorname{Exp}(\lambda)\) & \([0,\infty)\) &
  \(1/\lambda\) & \(1/\lambda^2\)\\
\(\operatorname{Poisson}(m)\) & \(0,1,\ldots\) & \(m\) & \(m\)\\
\(\operatorname{Gamma}(\alpha,\beta)\) & \([0,\infty)\) &
  \(\alpha/\beta\) & \(\alpha/\beta^2\)\\
\(\operatorname{Bernoulli}(p)\) & \(0,1\) & \(p\) & \(p(1-p)\)\\
\(\operatorname{Binomial}(n,p)\) & \(0,\ldots,n\) &
  \(np\) & \(np(1-p)\)\\
\(\operatorname{Geometric}(p)\) & \(1,2,\ldots\) &
  \(1/p\) & \((1-p)/p^2\)\\
\(\mathcal N(\mu,\sigma^2)\) & \(\Real\) & \(\mu\) & \(\sigma^2\)\\
\(t_\nu\) & \(\Real\) &
  \(0\;(\nu>1)\) & \(\nu/(\nu-2)\;(\nu>2)\)\\
\bottomrule
\end{tabularx}
\end{center}
Here \(\lambda,\alpha,\beta,m,\nu>0\), \(\sigma>0\), and \(n\) is a
nonnegative integer. The Bernoulli and Binomial parameters allow
\(0\le p\le1\); the Geometric requires \(0<p\le1\).
Gamma uses shape and rate. Geometric counts trials including the first
success. Student moments exist only under the conditions shown.

\begin{center}
\begin{tabularx}{\textwidth}{@{}L{0.38\textwidth}Y@{}}
\toprule
\textbf{Connection} & \textbf{Required mechanism or condition}\\
\midrule
Exponential to Gamma & Sum independent waiting times with a common rate.\\
Exponential to Poisson & Count arrivals generated by independent exponential
  interarrival times.\\
Bernoulli to Binomial & Sum a fixed number of independent trials with the
  same success probability.\\
Binomial to Poisson & Take \(n\to\infty\), \(p\to0\), with \(np=m\).\\
Geometric to Exponential & Rescale trial duration to \(\Delta\), set
  \(p=\lambda\Delta\), and let \(\Delta\downarrow0\).\\
Poisson to negative binomial & Mix the Poisson rate with a Gamma distribution.\\
Gaussian to Student & Mix the Gaussian precision with
  \(\operatorname{Gamma}(\nu/2,\nu/2)\).\\
\bottomrule
\end{tabularx}
\end{center}
